\documentclass[11pt]{article}
\usepackage[margin=1in]{geometry}
\usepackage{enumitem}
% =============================================================================
% Preambles/header.tex — shared LaTeX/Beamer preamble for this project
%
% Use from a Beamer lecture by prepending:
%     \input{header}
% and compiling with TEXINPUTS=../Preambles:$TEXINPUTS (the /compile-latex
% skill does this automatically).
%
% PALETTE CONTRACT. Color names defined here MUST match the names in
% Quarto/theme-template.scss so Beamer and Quarto renderings use the same
% palette. The scripts/check-palette-sync.sh script enforces this.
%
% To customize: edit the HEX values below AND the matching $variable in
% Quarto/theme-template.scss, then run scripts/check-palette-sync.sh.
% =============================================================================

% -----------------------------------------------------------------------------
% Engine detection + encoding
%
% Our /compile-latex skill uses XeLaTeX, where inputenc is unnecessary and
% can produce encoding warnings. Only load inputenc under pdfTeX.
% -----------------------------------------------------------------------------
\usepackage{iftex}
\ifPDFTeX
  \usepackage[utf8]{inputenc}
\fi

\usepackage{amsmath, amssymb, amsthm}
\usepackage{booktabs}
\usepackage{graphicx}

% -----------------------------------------------------------------------------
% PALETTE — mirrors Quarto/theme-template.scss variable names.
% Load xcolor and define colors BEFORE anything that references them
% (notably hyperref, hypersetup, and the Beamer color assignments).
% -----------------------------------------------------------------------------
\usepackage{xcolor}

\definecolor{primary-blue}{HTML}{012169}      % headings, accents
\definecolor{primary-gold}{HTML}{B9975B}      % emphasis, borders
\definecolor{highlight-yellow}{HTML}{F2A900}  % markers, alerts
\definecolor{light-bg}{HTML}{E8EDF5}          % title / section backgrounds
\definecolor{jet}{HTML}{1A1A1A}               % body text

% Semantic colors (used in diagrams and annotations)
\definecolor{positive}{HTML}{15803D}          % good / observed / identified
\definecolor{negative}{HTML}{B91C1C}          % bad / problematic / bias
\definecolor{neutral}{HTML}{525252}           % reference / context

% Highlight accents (match .hi-* classes in the SCSS)
\definecolor{hi-slate}{HTML}{314F4F}
\definecolor{hi-green}{HTML}{2E8B57}
\definecolor{hi-red}{HTML}{C41E3A}

% Semantic aliases so skill templates and snippets can write
% \textcolor{accent}{...} without hard-coding "primary-blue".
\colorlet{accent}{primary-blue}
\colorlet{accent2}{primary-gold}

% -----------------------------------------------------------------------------
% Hyperref — loaded AFTER xcolor + palette so link colors resolve.
% -----------------------------------------------------------------------------
\usepackage{hyperref}
\hypersetup{colorlinks=true, linkcolor=primary-blue, urlcolor=primary-blue,
            citecolor=primary-blue, pdfborder={0 0 0}}

% -----------------------------------------------------------------------------
% Course code — set once per deck (after \input{header}, before \begin{document})
% via \coursecode{CS401}. Shown in the footer on every slide so a deck's course
% is identifiable without opening the title page. Empty by default.
% -----------------------------------------------------------------------------
\makeatletter
\newcommand{\coursecode}[1]{\gdef\@coursecodetext{#1}}
\gdef\@coursecodetext{}
\makeatother

% -----------------------------------------------------------------------------
% Beamer theme assignments (only apply under Beamer).
%
% We detect Beamer via \@ifclassloaded{beamer}, which is the documented,
% non-internal way. This must be wrapped in \makeatletter/\makeatother
% because \@ifclassloaded contains an @-catcode character.
% -----------------------------------------------------------------------------
\makeatletter
\@ifclassloaded{beamer}{%
  \usefonttheme{professionalfonts}
  \setbeamercolor{structure}{fg=primary-blue}
  \setbeamercolor{frametitle}{fg=primary-blue}
  \setbeamercolor{title}{fg=primary-blue}
  \setbeamercolor{subtitle}{fg=primary-gold}
  \setbeamercolor{author}{fg=primary-blue}
  \setbeamercolor{institute}{fg=neutral}
  \setbeamercolor{date}{fg=primary-gold}
  \setbeamercolor{itemize item}{fg=primary-blue}
  \setbeamercolor{itemize subitem}{fg=primary-gold}
  \setbeamercolor{alerted text}{fg=primary-gold}
  \setbeamercolor{block title}{fg=white, bg=primary-blue}
  \setbeamercolor{block body}{bg=light-bg}
  % Semantic block variants: block=definition/concept (blue), exampleblock=
  % worked example (green/positive), alertblock=key takeaway or caution (gold).
  % Keep to <=2 colored boxes per slide across ALL three variants (INV-7).
  \setbeamercolor{block title example}{fg=white, bg=positive}
  \setbeamercolor{block body example}{bg=light-bg}
  \setbeamercolor{block title alerted}{fg=white, bg=primary-gold}
  \setbeamercolor{block body alerted}{bg=light-bg}

  % Minimal footer: course code (left) + page X / N (right), no navigation clutter
  \setbeamertemplate{navigation symbols}{}
  \setbeamertemplate{footline}{%
    \color{neutral}\footnotesize\hspace{0.7em}\@coursecodetext\hfill
    \insertframenumber/\inserttotalframenumber\hspace{0.7em}\vspace{0.3em}}
}{}
\makeatother

% -----------------------------------------------------------------------------
% TikZ setup — libraries the snippet gallery uses
% -----------------------------------------------------------------------------
\usepackage{tikz}
\usetikzlibrary{arrows.meta, positioning, calc, decorations.pathreplacing,
                fit, shapes.geometric, backgrounds}

% Shared TikZ styles — snippets and hand-written diagrams can pick these up
% via `\begin{tikzpicture}[every node/.style={...}]` or by naming specific
% styles (e.g., `\node[dag-node] (X) at (0,0) {X};`).
\tikzset{
  % Node styles for causal DAGs and similar
  dag-node/.style = {
    draw=primary-blue, fill=light-bg, rounded corners=2pt,
    minimum width=1.2cm, minimum height=0.9cm, align=center, font=\small
  },
  decision-node/.style = {
    draw=primary-gold, fill=white, diamond, aspect=1.8,
    minimum width=1.8cm, minimum height=1.2cm, align=center, font=\small,
    inner sep=1pt
  },
  % Edge styles — observed vs counterfactual
  observed-edge/.style = {-{Latex[length=2.5mm]}, thick, draw=jet},
  counterfactual-edge/.style = {-{Latex[length=2.5mm]}, thick, draw=neutral, dashed},
  confound-edge/.style = {-{Latex[length=2.5mm]}, thick, draw=negative, dashed},
  % Data-point styles for plots
  observed-dot/.style = {fill=primary-blue, draw=primary-blue, circle, inner sep=1.2pt},
  counterfactual-dot/.style = {fill=white, draw=neutral, circle, inner sep=1.2pt}
}

% -----------------------------------------------------------------------------
% Convenience macros
% -----------------------------------------------------------------------------
\newcommand{\muted}[1]{\textcolor{neutral}{#1}}
\newcommand{\key}[1]{\textcolor{primary-gold}{\textbf{#1}}}
\newcommand{\good}[1]{\textcolor{positive}{#1}}
\newcommand{\bad}[1]{\textcolor{negative}{#1}}

% Standout frame macro: full-bleed dark background for section transitions.
% Beamer-only (uses \begin{frame}/\setbeamercolor/\usebeamerfont) -- do not
% call from an article-class file (e.g. Notes/<CODE>/*-notes.tex). \muted,
% \key, \good, \bad, and \coursecode above are plain xcolor/gdef and are
% safe to use from either class; verified when Notes/article-class support
% was added.
% Expand: \transitionslide{Your transition title}
\newcommand{\transitionslide}[1]{%
  \begingroup
  \setbeamercolor{background canvas}{bg=primary-blue}
  \begin{frame}[plain]
    \centering
    \vfill
    {\usebeamerfont{title}\color{white}\Large #1\par}
    \vfill
  \end{frame}
  \endgroup
}

% Section divider frame (Beamer-only), an alternative to \transitionslide with a
% darker two-line style: near-black (jet) background, a small white label line
% above a larger gold title, and a thin gold rule along the bottom edge.
% Expand: \sectiondivider{Section 2}{Pointers}
\newcommand{\sectiondivider}[2]{%
  \begingroup
  \setbeamercolor{background canvas}{bg=jet}
  \begin{frame}[plain]
    \vspace*{3.0cm}
    {\color{white}\ttfamily\large #1\par}
    \vspace{0.5cm}
    {\color{primary-gold}\LARGE #2\par}
    \begin{tikzpicture}[remember picture, overlay]
      \draw[primary-gold, line width=2.5pt]
        ($(current page.south west)+(0,0.1)$) -- ($(current page.south east)+(0,0.1)$);
    \end{tikzpicture}
  \end{frame}
  \endgroup
}

% -----------------------------------------------------------------------------
% Channel watermark — "Concepts That Click" (YouTube).
%
% Article-class files (CompetitiveExam Q/A sets, Notes, Assignments) get a
% light diagonal draft watermark on every page plus a centered footer naming
% the channel. Beamer decks get a subtle TikZ background watermark instead
% (the footline already carries the course code). Centralized here so every
% MCQ PDF is branded without per-file edits.
% -----------------------------------------------------------------------------
\makeatletter
\@ifclassloaded{article}{%
  \usepackage{draftwatermark}
  \SetWatermarkText{Concepts That Click}
  \SetWatermarkScale{1}
  \SetWatermarkFontSize{2cm}
  \SetWatermarkLightness{0.86}
  \SetWatermarkAngle{45}
  \usepackage{fancyhdr}
  \pagestyle{fancy}
  \fancyhf{}
  \fancyfoot[C]{\footnotesize\textcolor{neutral}{Concepts That Click~|~YouTube: \href{https://www.youtube.com/@concepts.thatclick}{@concepts.thatclick}}}
  \renewcommand{\headrulewidth}{0pt}
  \renewcommand{\footrulewidth}{0pt}
  \fancypagestyle{plain}{%
    \fancyhf{}%
    \fancyfoot[C]{\footnotesize\textcolor{neutral}{Concepts That Click~|~YouTube: \href{https://www.youtube.com/@concepts.thatclick}{@concepts.thatclick}}}%
    \renewcommand{\headrulewidth}{0pt}%
    \renewcommand{\footrulewidth}{0pt}%
  }
}{}
\@ifclassloaded{beamer}{%
  \setbeamertemplate{background}{%
    \begin{tikzpicture}[remember picture, overlay]
      \node[rotate=45, opacity=0.055, align=center]
        at (current page.center)
        {\Huge\textcolor{neutral}{Concepts That Click}};
    \end{tikzpicture}%
  }
}{}
\makeatother 
\coursecode{JKSSB}
\renewcommand{\thesection}{56.\arabic{section}}
\newtheorem{example}{Example}[section]
\newenvironment{definitionbox}[1]{%
  \par\noindent\textbf{Definition (#1).}\ \itshape
}{\par\vspace{0.3em}}
\title{Cybersecurity, IT Act and Security Controls --- Detailed Lecture Notes}
\author{JKSSB Computer Awareness}
\date{\today}

\begin{document}
\maketitle

\section{Why security, and the CIA triad}
A working computer holds records worth protecting: student marks, account
details, personal files, and official correspondence on the office desktop.
\textbf{Cybersecurity} is the practice of protecting computers, networks,
and data against theft, damage, unauthorised access, and misuse. Three
goals organise the whole topic, remembered as the \textbf{CIA triad}:
\textbf{confidentiality} (only authorised persons can read the data),
\textbf{integrity} (the data is not altered by unauthorised means), and
\textbf{availability} (the system stays usable when needed). Every threat
in this lesson attacks one of these goals, and every control defends one.

\begin{definitionbox}{Cybersecurity}
The protection of computer systems, networks, and data against theft,
damage, disruption, and unauthorised access, so that confidentiality,
integrity, and availability are preserved.
\end{definitionbox}

\begin{center}
\begin{tabular}{p{0.24\textwidth}p{0.66\textwidth}}
\toprule
\textbf{Term} & \textbf{Meaning in this lesson} \\
\midrule
Malware & Umbrella term for malicious software written to harm or exploit. \\
Virus & Malware that attaches to a host file and needs human action to spread. \\
Worm & Standalone self-copying malware that spreads over a network unaided. \\
Trojan horse & Malware disguised as useful software; opens a backdoor; does not replicate. \\
Ransomware & Malware that locks or encrypts files and demands payment. \\
Encryption & Coding data so only authorised persons can read it. \\
SSL / TLS & Secure Sockets Layer / Transport Layer Security; secure transmission protocols. \\
Firewall & Barrier filtering network traffic against rules. \\
Quarantine & Isolation of a suspicious file so it cannot run. \\
MFA & Multi-Factor Authentication; two or more identity checks at login. \\
\bottomrule
\end{tabular}
\end{center}

The rest of this lesson uses these terms in one consistent sense.
\textbf{Malware} is always the general term, and virus, worm, Trojan, and
ransomware are its specific types. \textbf{Encryption} always means hiding
content (confidentiality), while \textbf{digital signature} always means
proving the sender and intactness (authenticity and integrity). The exam
tests exactly these pairings, so the wording here is kept fixed.

\section{The Information Technology Act, 2000}
The \textbf{Information Technology Act, 2000} is India's law for the
electronic world. Its scope has three parts: it gives \textbf{legal
validity to electronic records and electronic signatures}, so that online
transactions and electronic governance stand on a legal footing; it enables
\textbf{electronic governance} (official records, notices, and transactions
in electronic form); and it defines \textbf{cyber contraventions and
offences}, such as unauthorised access to a computer, damage to computer
systems and data, and fraud or cheating committed through a computer.

\begin{definitionbox}{Information Technology Act, 2000}
India's legislation of the year 2000 that recognises electronic records
and electronic signatures in law, supports electronic governance and
electronic commerce, and provides for cyber contraventions and offences.
\end{definitionbox}

Two exam points follow. First, the year is \textbf{2000}; later changes to
the law were amendments to the 2000 Act, not a fresh Act, so an option
calling the Act itself ``the IT Act of 2008'' is wrong. Second, the scope
words are \textbf{electronic records, electronic governance, and
cybercrime}: an option that limits the Act to hardware prices, typing
tests, or office furniture has left the scope entirely. No section numbers
are needed for this lesson; the exam fact is the Act's year and its three
branches of coverage.

\section{Malware: the umbrella and its members}
\textbf{Malware}, short for malicious software, is any program written to
harm a computer, steal data, or disrupt work. It is the umbrella term: the
question stem may say ``malware'' when it means the general class, and name
virus, worm, Trojan, or ransomware when it means one member. Two further
nuisance programs belong nearby: \textbf{spyware}, which secretly watches
the user's activity and reports it, and \textbf{adware}, which displays
unwanted advertisements.

\begin{definitionbox}{Malware}
Malicious software: any program designed to damage, disrupt, or gain
unauthorised access to a computer system. Virus, worm, Trojan horse, and
ransomware are types of malware.
\end{definitionbox}

A \textbf{virus} attaches itself to a host file or program and spreads when
that file is run or shared. It needs human action --- opening an infected
attachment, running an infected file, or carrying it on a pen drive --- and
once it runs it can corrupt, delete, or alter files. A \textbf{worm}, by
contrast, is standalone: it copies itself and spreads across a network
without any host file and without waiting for a user to click. Worms
consume bandwidth and system resources, so their signature effect is a
slowed or crashed network. A \textbf{Trojan horse} pretends to be useful
software --- a free utility, a game, a handy tool --- but hides a harmful
purpose inside. It does not replicate by itself; instead it opens a
\textbf{backdoor} through which an attacker can steal passwords, spy on the
user, or control the machine. \textbf{Ransomware} encrypts the victim's
files or locks the screen and demands payment (a ransom) for release.
Paying does not guarantee recovery; the reliable defence is a clean
\textbf{backup} kept separate from the machine.

\begin{example}
An item describes a program that ``looks like a useful download, does not
copy itself, but lets an attacker read passwords.'' Trace the three
contrasts: it needs no host file to be described (so not a virus), it does
not self-spread (so not a worm), and it demands no payment (so not
ransomware). The false-useful disguise plus the backdoor is the Trojan
horse.
\end{example}

\section{Antivirus, quarantine, and patching}
\textbf{Antivirus} software detects, blocks, and removes malware. It works
by recognising known malware signatures and by flagging suspicious
behaviour, and it must be updated regularly because new malware appears
every day. When the antivirus meets a file it cannot yet judge, it places
it in \textbf{quarantine}: an isolated holding area where the file cannot
run or spread while the user or administrator decides its fate. Quarantine
is isolation, not deletion and not a backup --- the file is still present
but harmless while isolated.

\begin{definitionbox}{Quarantine}
The isolation of a suspicious file by security software so that it cannot
execute or spread, pending inspection, removal, or restoration.
\end{definitionbox}

\textbf{Patching} is the companion hygiene step: installing the
vendor-supplied updates that fix known security flaws in the operating
system and applications. An unpatched machine remains open to attacks that
the vendor has already fixed, so ``install updates promptly'' is the exam's
standard correct advice. Antivirus cleans what is already inside,
quarantine holds what is suspect, and patching closes the holes through
which the next attack would enter.

\section{Phishing and anti-phishing}
\textbf{Phishing} is not malware at all but deception: a fake message ---
mail, SMS, or website --- that impersonates a trusted sender (a bank, an
office, a delivery firm) to steal passwords, OTPs, or money. Its warning
signs are urgency and threat (``your account will be blocked''), unexpected
attachments, links whose real destination does not match the claimed
sender, and any request for a password or OTP. The defence, called
\textbf{anti-phishing} practice, is procedural: verify the sender through a
separate channel, hover over or long-press links to inspect them before
clicking, never share an OTP or password by message, and report the message
to the service provider or office.

\begin{definitionbox}{Phishing}
A fraudulent message impersonating a trusted sender, designed to trick the
recipient into revealing passwords, OTPs, or financial details.
\end{definitionbox}

\begin{example}
An item shows a mail reading ``Dear customer, your account is blocked;
click here and enter your OTP to restore it.'' The correct response is to
treat it as phishing: do not click, do not enter the OTP, and verify with
the bank through its known number or site. A genuine bank or office never
asks for a password or OTP by mail or message, so any option advising
compliance is the trap.
\end{example}

\section{Encryption, digital signatures, and SSL/TLS}
\textbf{Encryption} converts readable data (plaintext) into coded form
(ciphertext) so that only authorised persons holding the key can read it.
Its goal is confidentiality. A \textbf{digital signature}, by contrast,
proves who sent a message and that it was not altered in transit; its goals
are authenticity and integrity. The pair must never be swapped: a signature
does not hide the message (anyone can still read it), and encryption by
itself does not prove who sent it. This is the same distinction tested in
the email lesson (TRAP-15), and it returns here because ransomware's
``encryption'' step uses the same idea against the victim.

\textbf{SSL} (Secure Sockets Layer) and its successor \textbf{TLS}
(Transport Layer Security) provide a secure, encrypted channel over a
network. They protect data \textbf{in transit}: mail submission, web
logins, and online forms, signalled by the padlock in the browser address
bar. The lesson's anchor fact, already seen in WO 2026 Q76, is that the
secure email transmission protocol is \textbf{SSL/TLS}. Note the limit:
SSL/TLS protect the journey, not the endpoints --- they do not clean an
infected file and do not replace antivirus.

\begin{definitionbox}{Encryption}
The conversion of plaintext into ciphertext with a key, so that only
authorised holders of the key can recover the original content.
\end{definitionbox}

\section{Firewall, MFA, and the matching map}
A \textbf{firewall} is a barrier, implemented in hardware or software (or
both), that filters incoming and outgoing network traffic against a set of
rules. It blocks unauthorised access while letting permitted traffic pass,
like a gatekeeper at the office network door checking every packet against
the list. Its partner inside the machine is the antivirus: the firewall
filters traffic at the boundary, while the antivirus cleans malware already
inside. One does not replace the other.

\textbf{MFA} (Multi-Factor Authentication) protects logins by demanding two
or more independent checks --- for example, a password plus an OTP, or a
password plus a fingerprint. It closes the door against stolen passwords:
even a leaked password is useless without the second factor. Strong,
distinct passwords remain the first factor; MFA adds factors rather than
replacing them.

\begin{center}
\begin{tabular}{p{0.28\textwidth}p{0.62\textwidth}}
\toprule
\textbf{Control} & \textbf{Job (the stem verb to watch)} \\
\midrule
Firewall & Filters network traffic; blocks unauthorised access. \\
Antivirus & Detects and removes malware on the machine. \\
Quarantine & Isolates a suspicious file so it cannot run. \\
Encryption / SSL-TLS & Hides content; protects data in transit. \\
Digital signature & Proves the sender and that content is intact. \\
Patching & Fixes known security flaws by installing updates. \\
MFA & Adds extra login checks beyond the password. \\
Backup & Recovers data after ransomware, deletion, or failure. \\
\bottomrule
\end{tabular}
\end{center}

The matching technique is mechanical: read the verb in the stem.
``Filters traffic'' points to the firewall; ``cleans malware'' to the
antivirus; ``isolates'' to quarantine; ``hides content'' to encryption;
``proves sender'' to digital signature; ``fixes flaws'' to patching;
``extra login check'' to MFA; ``recovers after loss'' to backup.

\section{Exam retrieval and common confusions}
Seven distinctions prevent most errors on this topic.
\begin{enumerate}
  \item The IT Act is of \textbf{2000} and covers electronic records,
    e-governance, and cybercrime; later changes were amendments.
  \item Malware is the \textbf{umbrella}; virus, worm, Trojan, and
    ransomware are specific types under it.
  \item Virus needs a \textbf{host file plus human action}; worm is
    standalone and \textbf{self-spreading}; Trojan \textbf{does not
    replicate}; ransomware \textbf{locks files and demands payment}.
  \item A digital signature \textbf{proves} (authenticity/integrity); it
    does not \textbf{hide} --- encryption provides confidentiality
    (TRAP-15).
  \item \textbf{SSL/TLS} protect data \textbf{in transit}; they do not
    clean infected files.
  \item A \textbf{firewall} filters \textbf{traffic} at the boundary;
    antivirus cleans \textbf{malware} inside; quarantine
    \textbf{isolates} without deleting.
  \item \textbf{MFA} adds login factors; \textbf{patching} fixes software
    flaws; \textbf{backup} is the recovery answer to ransomware.
\end{enumerate}

\begin{example}
An item asks which control ``proves that a message came from the claimed
sender and was not altered.'' The verb ``proves'' selects the digital
signature, not encryption (which hides) and not the firewall (which
filters). If the stem instead says ``hides the content from
eavesdroppers,'' the same map selects encryption; if it says ``blocks
unauthorised packets,'' it selects the firewall.
\end{example}

\subsection*{Exercises}
\begin{enumerate}[label=\arabic*.]
  \item State the year and the three branches of scope of the Information
    Technology Act covered in this lesson.
  \item Define malware and place virus, worm, Trojan horse, and ransomware
    under it in one sentence each.
  \item Distinguish a virus from a worm on the two points of host file and
    human action.
  \item Why is a Trojan horse dangerous even though it does not replicate?
  \item Explain why backup, rather than payment, is the correct response to
    ransomware.
  \item Contrast encryption with a digital signature: which goal does each
    serve?
  \item Using the matching map, assign one control to each verb: filters,
    isolates, hides, proves, fixes.
\end{enumerate}

\subsection*{Solutions}
\begingroup\small
\begin{enumerate}[label=\arabic*.]
  \item The Information Technology Act is of \textbf{2000}. It gives legal
    validity to electronic records and electronic signatures, enables
    electronic governance, and defines cyber contraventions and offences.
  \item Malware is malicious software written to harm or exploit. A virus
    attaches to a host file and needs human action; a worm self-copies
    across networks unaided; a Trojan disguises itself as useful software
    and opens a backdoor without replicating; ransomware locks or encrypts
    files and demands payment.
  \item A virus needs a host file and a human act (opening, running,
    sharing) to spread; a worm is standalone, copies itself, and spreads
    over the network without a host file or user action.
  \item The Trojan's disguise persuades the user to install it, and the
    hidden backdoor then lets the attacker steal passwords, spy, or
    control the machine --- damage that needs no replication.
  \item Payment does not guarantee recovery and rewards the attacker,
    while a clean separate backup restores the data without negotiation;
    hence backup is the reliable defence.
  \item Encryption serves \textbf{confidentiality} (hides content);
    a digital signature serves \textbf{authenticity and integrity}
    (proves sender and intactness).
  \item Filters --- firewall; isolates --- quarantine; hides ---
    encryption (SSL/TLS in transit); proves --- digital signature;
    fixes --- patching (installing updates).
\end{enumerate}
\endgroup
\end{document}
